\abstract{%
We develop a generalized hierarchical Bayesian method for multiple-changepoint segmentation of ordered data observed on regular or irregular designs. For a regular exponential-family observation model with a proper segmentwise conjugate prior and finite normalizing constants, cumulative sufficient statistics determine the marginal likelihood and posterior moments of every candidate segment. These segment marginal likelihoods enter dynamic-programming recursions for the marginal likelihood conditional on the number of segments, the posterior distribution of the segment count, marginal changepoint probabilities, and posterior moments of the piecewise-constant signal. A separate max-sum recursion with backtracking returns the joint maximum-a-posteriori partition. Thus, the partition calculation is common across observation families once the segment marginal likelihoods are specified, while the statistical meaning of each segment remains family-specific.

The hierarchical formulation includes design-dependent partition priors for irregular coordinates, common-changepoint models for multiple conditionally independent sequences, known-group models, and a finite latent-group template model. The latent-group algorithm uses a Jensen minorization and exact responsibility-weighted max-sum updates for the stated finite objective; because its component scores are not normalized sampling densities, it is not interpreted as a finite-mixture likelihood. When a nonconjugate segment integral is replaced by a deterministic approximation, a uniform error bound on reachable segment log marginal likelihoods implies explicit bounds on partition marginal likelihoods and posterior odds.

Simulation studies report boundary-posterior calibration for Gaussian data
(expected calibration error approximately $0.010$, Brier score approximately
$0.011$), applications of common partition recursions to Gaussian, Poisson,
Binomial, and Beta-valued responses, and 96\% hard allocation accuracy in one
two-group template experiment. A separate 400-dataset stress study evaluates
template recovery and objective monotonicity. Observed Gaussian computation
times range from $0.0455$ seconds at $n=50$ to $1.0585$ seconds at $n=800$
for $k_{\max}=20$. Analyses of well-log, array-CGH, financial-return, and
methylation data illustrate model use, while the lack of independently scored
changepoints limits conclusions about external segmentation accuracy. Additional
simulation and comparison protocols are specified as planned work and do not
contribute numerical findings to this report.}